\documentclass[a4paper,11pt]{article}
\usepackage[margin=24mm]{geometry}
\usepackage{xeCJK}
\setCJKmainfont[BoldFont=HaranoAjiMincho-Bold.otf,ItalicFont=HaranoAjiMincho-Regular.otf]{HaranoAjiMincho-Regular.otf}
\setCJKsansfont{HaranoAjiGothic-Medium.otf}
\setCJKmonofont{HaranoAjiGothic-Medium.otf}
\usepackage{amsmath,amssymb,amsthm,mathtools}
\usepackage{booktabs,longtable,array,enumitem}
\usepackage{xcolor}
\usepackage[unicode,colorlinks=true,linkcolor=blue!50!black,urlcolor=blue!50!black,citecolor=blue!50!black]{hyperref}
\usepackage{xurl}
\usepackage{bookmark}
\numberwithin{equation}{section}
\newtheorem{definition}{定義}[section]
\newtheorem{proposition}[definition]{命題}
\newtheorem{example}[definition]{例}
\renewcommand{\proofname}{証明}
\renewcommand{\contentsname}{目次}
\renewcommand{\abstractname}{概要}
\renewcommand{\refname}{参考文献}
\newcommand{\N}{\mathbb{N}_0}
\newcommand{\R}{\mathbb{R}}
\newcommand{\ind}{\mathbf{1}}
\newcommand{\emptycell}{\mathsf{empty}}
\newcommand{\garb}{\mathsf{garbage}}
\newcommand{\unknown}{\mathsf{unknown}}
\newcommand{\none}{\mathsf{none}}
\newcommand{\code}[1]{\nolinkurl{#1}}
\newcommand{\src}[2]{\href{https://github.com/shhchan/puyo-ai-dev-platform/blob/f8394d77578ecb66e15802fa502d8e1668223006/#1}{\code{#2}}}
\newcommand{\srcja}[2]{\href{https://github.com/shhchan/puyo-ai-dev-platform/blob/f8394d77578ecb66e15802fa502d8e1668223006/#1}{#2}}
\setlist{itemsep=2pt,topsep=4pt}
\setlength{\emergencystretch}{3em}
\title{ぷよぷよの状態遷移・連鎖探索・戦術強化学習\\\large nextgen\_tactic\_manager の数理的定式化}
\author{puyo-ai-dev-platform / PUYO-282}
\date{2026 年 10 月 13 日\quad 第 1 版}
\begin{document}
\maketitle
\begin{abstract}
有限盤面上の消去ゲームを出発点とし，公開情報に制限された候補探索と，有限個の戦術を選択する学習方策の合成を定義する．具体的な配置は固定探索器に，戦術の継続と切替は学習可能な選択器に担当させる．ゲームの得点，探索用評価値，対戦学習の報酬を区別し，時刻，未観測情報，探索打切り，実行失敗を明示する．PR \#188 の実装事実と，PUYO-234 の構想を具体化する将来仕様を区別する．
\end{abstract}
\setcounter{tocdepth}{1}
\tableofcontents
\clearpage
\section{対象・定義の読み方}\label{sec:scope}
本稿の「ぷよぷよ」は本リポジトリの 4 色対戦規則を指す．製品や大会規則全般との同一性は主張しない．実装の基準は
\href{https://github.com/shhchan/puyo-ai-dev-platform/pull/188}{PR \#188} の
\code{f8394d77578ecb66e15802fa502d8e1668223006} である．
参照を可変なブランチ名に依存させない．構想の根拠は同じ版の
\srcja{docs/development/puyo-234-next-generation-model-design.md}{PUYO-234 設計書}である．

\textbf{実装記述}はこの版の動作，\textbf{抽象化}はその動作を扱うための数学的な型や写像，
\textbf{提案}は今後の学習器に課す仕様を表す．現行の選択器は
\code{RuleTacticSelector} であり，このモデル専用の学習済み戦術方策が接続済みとは扱わない．
既存の別モデル用 PPO と，本稿の 6 戦術 PPO を区別する．

記号は本文で初めて使用する直前に導入する．関数を設定として与える場合も定義域・値域・役割を明示し，
未実装の関数を既存関数と偽らない．一般の集合・確率・線形代数の記法は次の規約を用いる．
\begin{definition}[基礎記法]\label{def:notation}
$\N=\{0,1,2,\ldots\}$ は非負整数，$\R$ は実数，$\R_{\geq0}$ は非負実数の集合である．
$|A|$ は有限集合 $A$ の要素数，$A\times B$ は直積，$B^A$ は $A$ から $B$ への全写像の集合，
$A^*$ は $A$ の有限列全体である．$\varnothing$ は空集合である．
条件 $P$ の指示値 $\ind\{P\}$ は成立時に $1$，それ以外で $0$ とする．
$[v]_+=\max(0,v)$，$\operatorname{clip}(v;l,u)=\min(u,\max(l,v))$ とする．
$\lfloor v\rfloor$ は床関数，$\lceil v\rceil$ は天井関数，$\log$ は自然対数，$\exp$ はその逆関数である．
$\sum$ と $\prod$ は添字集合上の和と積であり，空和は $0$，空積は $1$ とする．
$\none$ は「該当値なし」を表す特別な記号で，数値 $0$ とは異なる．

有限または可算集合 $A$ に対する $\Delta(A)$ は，非負で総和 $1$ の確率質量関数全体とする．
$X\sim p$ は確率変数 $X$ が分布 $p$ に従うこと，
$\Pr$ は指定した実験の確率，$\mathbb E$ はその期待値を表す．
条件付き期待値は条件の確率が正の範囲で用いる．連続変数を持つ一般化では対応する確率測度と積分に置き換える．
$\arg\max$，$\arg\min$ は最大値・最小値を与える要素の集合を表し，
単一要素を必要とする箇所では必ず指定した全順序で同点を解消する．
辞書式順序は，左から最初に異なる成分の大小で列を比較する順序である．
$\circ$ は写像の合成である．
\end{definition}

\section{盤面，配置，連鎖，得点}\label{sec:game}
\subsection{盤面と重力}
\begin{definition}[盤面の型]\label{def:board}
色集合を $\mathcal C=\{\mathsf r,\mathsf b,\mathsf g,\mathsf y\}$，空セルを $\emptycell$，
おじゃまを $\garb$ とする．セル値集合は $\mathcal V=\mathcal C\cup\{\emptycell,\garb\}$ である．
幅 $W=6$，全高 $H=14$，消去対象の高さ $H_v=12$ を固定し，
$x$ を左から，$y$ を下から数える座標とする．盤面領域 $D$，可視領域 $D_v$，盤面集合 $\mathcal B$ を
\begin{equation}\label{eq:board}
 D=\{0,\ldots,5\}\times\{0,\ldots,13\},\quad
 D_v=\{0,\ldots,5\}\times\{0,\ldots,11\},\quad
 \mathcal B=\mathcal V^D
\end{equation}
と定める．盤面 $B\in\mathcal B$ に対して $B(x,y)$ はセル値である．
$y=12$ は落下可能だが消去判定の対象外，$y=13$ は一度置くと落下も消去もしない永続セルである．
領域外は占有済みの壁として衝突判定するが，壁を $D$ 内のセル値には含めない．
\end{definition}

列 $x$ に対して，$y=0,\ldots,12$ を昇順に走査して空以外を抽出した列を
$(v_{x,0},\ldots,v_{x,n_x-1})$ とする．$n_x$ はその長さである．
重力写像 $G:\mathcal B\to\mathcal B$ を
\begin{equation}\label{eq:gravity}
 G(B)(x,y)=\begin{cases}
 v_{x,y}&0\le y<n_x,\\
 \emptycell&n_x\le y\le12,\\
 B(x,13)&y=13
 \end{cases}
\end{equation}
で定める．$G(B)=B$ を満たす盤面を安定盤面という．
これは \src{src/core/field.py}{Field.drop_puyo} の規則であり，13 行全体の圧縮と永続行を分ける．

\subsection{消去と停止性}
$u=(x,y),v=(x',y')\in D_v$ に対し，$|x-x'|+|y-y'|=1$ のとき隣接するという．
色 $c\in\mathcal C$ のセルを頂点，隣接を辺とする無向グラフの極大連結成分を色成分という．
盤面 $B$ のうち要素数 $4$ 以上の色成分の集合を $\mathcal K(B)$ とする．
消える色セル集合 $V(B)$ と，それに直接隣接するおじゃま集合 $J(B)$ を
\begin{equation}\label{eq:erase}
 \begin{split}
 V(B)&=\bigcup_{K\in\mathcal K(B)}K,\\
 J(B)&=\{u\in D_v:B(u)=\garb,\ \exists v\in V(B)\text{ が }u\text{ に隣接}\}
 \end{split}
\end{equation}
と定める．消去写像 $E(B)$ は $V(B)\cup J(B)$ のセルを空にし，ほかを保持する写像である．
おじゃま同士の隣接による連鎖的な除去は行わない．

初期盤面 $B$ に対して $B^{(0)}=G(B)$ と置く．$\mathcal K(B^{(j)})\ne\varnothing$ の間だけ
$B^{(j+1)}=G(E(B^{(j)}))$ とする．最初に $\mathcal K(B^{(j)})=\varnothing$ となる添字を
連鎖数 $L(B)$ とし，解決後盤面を $\mathcal R(B)=B^{(L(B))}$ と定義する．
\begin{proposition}[連鎖解決の有限性]\label{prop:finite}
任意の $B\in\mathcal B$ について $L(B)\le21$ であり，$\mathcal R(B)$ は一意に定まる．
\end{proposition}
\begin{proof}
各消去で少なくとも $4$ 個の色セルが失われ，重力はセルを増やさない．初期の色セル数は
$|D|=84$ 以下なので，消去回数は $\lfloor84/4\rfloor=21$ 以下である．
各段の同時消去と重力は決定的なので一意である．この上限は緩い上限であり，達成可能性は主張しない．
\end{proof}

\subsection{組ぷよと配置行動}
組ぷよは順序付き色対 $p=(c_a,c_s)\in\mathcal C^2$ であり，第 1 成分を軸，第 2 成分を子とする．
向き集合を $\mathcal O=\{\mathsf U,\mathsf R,\mathsf D,\mathsf L\}$ とし，子の相対座標を
$\delta(\mathsf U)=(0,1)$，$\delta(\mathsf R)=(1,0)$，$\delta(\mathsf D)=(0,-1)$，
$\delta(\mathsf L)=(-1,0)$ と定める．$\delta_x,\delta_y$ はその各成分である．
配置行動集合 $\mathcal A$ は
\begin{equation}\label{eq:actions}
 \mathcal A=\{(x,o):x\in\{0,\ldots,5\},\ o\in\mathcal O,\ 0\le x+\delta_x(o)<6\},
 \qquad |\mathcal A|=22
\end{equation}
とする．$x$ の昇順，向きの順 $\mathsf U,\mathsf R,\mathsf D,\mathsf L$ で列挙した位置
$0,\ldots,21$ が実装の行動 ID である．例えば $(0,\mathsf U)$ は $0$，$(1,\mathsf L)$ は $6$ である．

衝突しない姿勢の集合 $\mathcal P(B)$ を，軸と子の両セルが $D$ 内かつ空である
$(x,y,o)$ 全体とする．同期配置における合法集合 $\mathcal A_{\rm geo}(B)$ は
$(x,12,o)\in\mathcal P(B)$ を満たす配置行動全体とする．その行動の着地点 $y_*(B,x,o)$ は，
$y=12$ から始め，$(x,y-1,o)\in\mathcal P(B)$ の間だけ $y$ を $1$ ずつ減らして得る値である．
軸・子を着地点に書き込む写像を $I(B,p,a)$ とし，同期配置解決を
\begin{equation}\label{eq:placement}
 B'=\mathcal R(I(B,p,a)),\qquad a\in\mathcal A_{\rm geo}(B)
\end{equation}
と定める．横向きの両セルが異なる高さに落ちる現象も $G$ が処理する．

\textbf{実装境界．} 式 \eqref{eq:placement} は配置単位の同期モデルである．
実時間では現在の軸位置，壁蹴り，入力リピート，接地期限から到達できない配置もある．
実時間制御状態を $\chi$（軸の座標・向き，押下状態，各入力と重力の次回 tick，接地回数・経過 tick）とする．
$\mathcal U$ を左右・下・左右回転キーの押下集合と解放集合の対全体とする．
設定済みの入力規則で $\chi$ から固定まで進む有限入力列 $u\in\mathcal U^*$ のうち，
指定期限までに行動 $a$ の配置に達するものが存在するとき，$a$ を到達可能と呼ぶ．
その集合を $\mathcal A_{\rm reach}(B,\chi)$ とする．
運用上は \src{puyo_env/action_planner.py}{action_planner} が証明できた到達可能集合を用い，
探索上限による未発見を物理的不可能の証明としない．
候補の根は $\mathcal A_{\rm geo}(B)\cap\mathcal A_{\rm reach}(B,\chi)$ に限定する．

\subsection{得点と全消し状態}
連鎖の第 $j$ 段（$1\le j\le L(B)$）について，消える色セル数を $n_j$，
同時に消える色の種類数を $d_j$，消える色成分集合を $\mathcal K_j=\mathcal K(B^{(j-1)})$ とする．
連鎖ボーナス関数 $b_{\rm ch}$，色数ボーナス関数 $b_{\rm col}$，連結ボーナス関数 $b_{\rm con}$ を
\begin{align}
 b_{\rm ch}(j)&=\begin{cases}0&j=1,\\8&j=2,\\16&j=3,\\32&j=4,\\32(\min(j,19)-3)&j\ge5,\end{cases}\label{eq:chainbonus}\\
 (b_{\rm col}(1),\ldots,b_{\rm col}(4))&=(0,3,6,12),\qquad
 b_{\rm con}(n)=\begin{cases}0&n=4,\\n-3&5\le n\le10,\\10&n\ge11\end{cases}\label{eq:otherbonus}
\end{align}
と定める．第 $j$ 段の素点 $s_j$ は
\begin{equation}\label{eq:score}
 s_j=10n_j\max\left(1,b_{\rm ch}(j)+b_{\rm col}(d_j)
                       +\sum_{K\in\mathcal K_j}b_{\rm con}(|K|)\right)
\end{equation}
である．消えたおじゃまは $n_j$ に加算しない．

全消しボーナス保留 bit を $b\in\{0,1\}$ とする．解決開始時の $b$ と盤面 $B$ に対して，
今回の連鎖得点 $S(B,b)$ と次の保留 bit $b'$ は
\begin{align}
 S(B,b)&=\sum_{j=1}^{L(B)}s_j+2100b\ind\{L(B)>0\},\label{eq:allclear-score}\\
 b'&=\begin{cases}
 b&L(B)=0,\\
 \ind\{\forall u\in D,\ \mathcal R(B)(u)=\emptycell\}&L(B)>0.
 \end{cases}\label{eq:allclear-bit}
\end{align}
全消しは隠れ行を含む $D$ 全体が空であることを要求し，ボーナスは次に成立した連鎖で消費する．
素点と操作による落下点の累積を総得点 $q\in\N$ とする．前回の連鎖終了時の総得点を $q_{\rm last}$ として保持する．
今回の連鎖終了時点の攻撃換算対象 $\Delta q$ は
\begin{equation}\label{eq:attackscore}
 \Delta q=\begin{cases}[q-q_{\rm last}]_+&L>0,\\0&L=0,\end{cases}
 \qquad q'_{\rm last}=\begin{cases}q&L>0,\\q_{\rm last}&L=0.\end{cases}
\end{equation}
ここで $L$ は今回の配置による連鎖数である．従って，前回連鎖以後の落下点も次の連鎖時に攻撃へ換算される．
同期配置では操作落下点がないため，その配置だけの素点と，対戦の攻撃差分を同一記号にしない．

\begin{example}[数値の区別]\label{ex:score}
同色 $4$ 個を消す $1$ 連鎖の素点は $40$ 点，その後さらに同色 $4$ 個を消す $2$ 段目は $320$ 点なので合計 $360$ 点となる．
保留ボーナスがなければ，換算率 $70$ 点でおじゃま $5$ 個と剰余 $10$ 点になる．
全消し成立時の盤面を空にするだけでは $2100$ 点を即時加算しない．
\end{example}

\section{実時間対戦と公開情報}\label{sec:versus}
\subsection{時刻とおじゃま台帳}
プレイヤー集合を $\mathcal I=\{0,1\}$ とし，$i\in\mathcal I$ の相手を $\bar i=1-i$ と書く．
試合 tick を $t\in\N$，tick rate を $f=60$（1 秒あたりの tick 数）とする．
おじゃま $1$ 個への換算点を $\kappa\in\{1,2,\ldots\}$（標準値 $70$），
攻撃換算の繰越点を $c_i\in\{0,\ldots,\kappa-1\}$ とする．
同一 tick の攻撃差分 $\Delta q_i$ に対し，生成個数 $g_i$ と新繰越点 $c_i'$ は
\begin{equation}\label{eq:carry}
 g_i=\left\lfloor\frac{\Delta q_i+c_i}{\kappa}\right\rfloor,
 \qquad c_i'=\Delta q_i+c_i-\kappa g_i.
\end{equation}
差分が $0$ の tick では攻撃を作らず，繰越点を保持する．

受信 packet は $(\mathrm{id},n,t_a,t_c,i_s)$ というレコードとする．
$\mathrm{id}$ は一意識別子，$n\in\N$ は残量，$t_a\in\N$ は到着可能 tick，
$t_c\in\N$ は生成 tick，$i_s\in\mathcal I$ は送信者である．
プレイヤー $i$ の未消費 packet の列を $Q_i$，その総量を $N_i$ とする．
到着前も含む受信 packet との相殺個数 $a_i$，相殺後生成量 $r_i$，両者同時相殺量 $a_*$，送信量 $o_i$ を
\begin{equation}\label{eq:cancel}
 N_i=\sum_{p\in Q_i}p.n,\quad a_i=\min(g_i,N_i),\quad
 r_i=g_i-a_i,\quad a_*=\min(r_0,r_1),\quad o_i=r_i-a_*
\end{equation}
と定める．$p.n$ は packet $p$ の残量である．$a_i$ 個を $(t_a,t_c)$ の昇順，同値なら元の列順で $Q_i$ から消費する．
$o_i>0$ なら相手に $(n=o_i,t_a=t+d_a,t_c=t,i_s=i)$ の新 packet を加える．
$d_a\in\N$ は固定攻撃遅延（標準値 $0$）である．

落下上限を $M\in\{1,2,\ldots\}$（標準値 $30$）とする．
受け手の配置・連鎖解決が終わった境界でのみ，到着済み残量から
$\min(M,\sum_{p\in Q_i}\ind\{p.t_a\le t\}p.n)$ 個の落下を試みる．
試行数を $n$ として $n=6h+r$，$h=\lfloor n/6\rfloor$，$0\le r<6$ と分解する．
全列へ $h$ 回，その後重複なく一様抽出した $r$ 列へ各 $1$ 個ずつ置く．
各列では $y=0,\ldots,11$ の最下の空セルへ挿入し，空セルがなければその試行を飛ばす．
台帳から消費するのは実際に配置した個数のみである．
到着済みでも操作中や連鎖中に割り込んで落下させない．残量は次の配置境界まで保持する．

\subsection{全状態と勝敗}
\begin{definition}[実時間状態と遷移]\label{def:state}
全状態集合 $\mathcal S$ を，両者の $B,p,\chi,q,q_{\rm last},b,c_i,Q_i$，
現在の phase，消去・落下の途中状態と残 tick，公開 event 履歴，
組ぷよ列とその読取位置，乱数生成器状態，試合終了 flag を含む有効レコード全体とする．
phase は操作中，連鎖演出中，生成待ち，おじゃま待ち，終了のいずれかである．
途中状態は現在の消去集合，連鎖段数，演出タイマーを含む．
再現に必要な有限長の内部カウンターも状態の一部とする．$s_t\in\mathcal S$ は tick $t$ の開始状態である．

規則設定全体を $\eta$ とする．$\eta$ は本節の定数，入力リピート・接地・壁蹴り規則，乱数生成方式を固定する．
$F_\eta:\mathcal S\times\mathcal U^2\to\mathcal S$ を，次の順に 1 tick を進める決定的写像として定義する．
(1) 両者の入力・重力または連鎖演出を進める，
(2) 両者の配置解決と窒息を調べる，
(3) 完了した連鎖の攻撃を式 \eqref{eq:carry}，\eqref{eq:cancel} で処理する，
(4) 試合終了過程でなければ到着済みおじゃまを配置境界で落とす，
(5) おじゃま待ち後に次の組を生成し，勝敗を更新する．
操作部分の具体的な $F_\eta$ は，固定 SHA の
\src{src/core/realtime.py}{RealtimeHeadlessSimulator.step} と
\src{puyo_env/realtime_versus.py}{RealtimeVersusMatch.step} の合成を採用する．
これは関数の実装による定義であり，全ての壁蹴り分岐を別の近似式へ置換するものではない．
\end{definition}

通常の窒息判定は生成境界またはおじゃま落下後の $B(2,11)\ne\emptycell$ である．
盤面最高点のみによる判定とは異なる．どちらかが窒息した後は新たな操作・生成・おじゃま落下を止め，
既に進行中の消去連鎖だけを完了させる．その処理が終わった集合を終端集合 $\mathcal S_T\subseteq\mathcal S$ とする．
$s\in\mathcal S_T$ におけるプレイヤー $i$ の結果 $z_i(s)\in\{-1,0,1\}$ は，
単独生存なら $1$，単独窒息なら $-1$，両者窒息なら総得点の大小により高い側 $1$，低い側 $-1$，同点 $0$ とする．
時間制限は $\mathcal S_T$ への遷移とはせず，後述の収集打切りとする．

組ぷよ生成方式は実験設定に含める．独立一様 4 色の理想化なら各対 $p\in\mathcal C^2$ の確率は $1/16$ である．
実装の \code{random} は seed 付き擬似乱数による同方式，\code{esports_tsu} は検証済み
$65536$ 配牌のうち選択した $128$ 組を周期反復する方式である．後者の試合内の対は独立同分布ではない．
配牌 ID の選択分布も実験設定として指定する．

\subsection{観測・履歴・非公開未来}
観測用セル集合を $\overline{\mathcal V}=\mathcal V\cup\{\unknown\}$ とし，$\unknown$ を未観測とする．
盤面の部分観測は $\bar B\in\overline{\mathcal V}^{D}$ であり，$\unknown\ne\emptycell$ である．
通信表現では空 $0$，通常色 $1,2,3,4$，おじゃま $5$，未観測 \code{None} を使い，上から下へ行を並べる．
14 行を送るとき行番号 $r$ と本稿の座標は $y=13-r$ で対応する．
行数が異なる観測は adapter の行配置を先に確定させ，この式を無条件には適用しない．

プレイヤー $i$ の公開観測 $o_t^i$ は，両者の部分盤面，公開済みの current/NEXT/NEXT2（各最大 $3$ 組），
公開 phase，公開 packet の量と到着・落下情報，繰越点，全消し状態，公開配置・消去等の event とする．
観測集合 $\mathcal O_{\rm pub}$ はこの型のレコード全体とする．
公開履歴 $h_t^i$ は時刻 $t$ までの $o^i$ と自分の入力・実行結果を時刻順に並べた有限列であり，
履歴集合を $\mathcal H$ とする．観測写像 $O_i$ は全状態・公開履歴からこの許可情報だけを抽出する手続きと定義する．
公開の配置履歴から一意に復元した隠れセルは観測へ追加できるが，非公開状態から補完しない．

公開部分盤面と整合する完全盤面の集合を
\begin{equation}\label{eq:completion}
 \mathcal B(\bar B)=\{B\in\mathcal B:\ \bar B(u)\ne\unknown\Rightarrow B(u)=\bar B(u)
 \text{ for every }u\in D\}
\end{equation}
とする．特定の補完盤面を使う予測はその補完を仮定した予測である．
全 $B\in\mathcal B(\bar B)$ に対する保証とは区別する．

探索・テンプレート選択・戦術抽選に用いる内部乱数列を $\xi$ とする．$\xi$ は環境の非公開 seed や配牌 ID から取得せず，
独立に指定する．候補・特徴・要求行動という意味的出力を返す写像を $\mathcal D$ とする．
計算の wall-clock 経過時間と scheduler の採用結果はこの写像の出力に含めない．情報境界の条件は
\begin{equation}\label{eq:privacy}
 h_t^i=h_t^{i\prime}\ \text{かつ同一の公開制御文脈と }\xi
 \quad\Longrightarrow\quad
 \mathcal D(h_t^i,\xi)=\mathcal D(h_t^{i\prime},\xi)
\end{equation}
である．右辺・左辺の非公開未来が異なっても，入力が同じなら出力は同じでなければならない．
公開観測の型検査だけではこの条件は証明できないので，producer 側で対照検証する．

\section{公開情報に基づく連鎖探索}\label{sec:search}
\subsection{探索問題とシナリオ}
探索時の単独状態を $x=(B,b)\in\mathcal X=\mathcal B\times\{0,1\}$ とする．
$B$ は安定盤面，$b$ は全消し保留 bit である．
遷移 $T(x,p,a)$ は式 \eqref{eq:placement} と \eqref{eq:allclear-bit} を適用した
$(x',L,S,d)$ を返す．$x'$ は次状態，$L$ は連鎖数，$S$ は今回の得点，
$d\in\{0,1\}$ は次の生成時の窒息を表す．非法行動では $T=\none$ とする．
この $T$ は相手の進行と受信 packet を含まない．対戦応答候補には別途実時間の検証が必要である．

根で公開されている組ぷよ列を $p_{0:k-1}=(p_0,\ldots,p_{k-1})$，$1\le k\le3$ とする．
探索深さ $D_s\ge1$，シナリオ数 $M_s\ge1$，幅 $w\ge1$，展開上限 $Q\ge1$，
根ごとの確保枠 $q_r\ge1$ を整数で与える．シナリオ $m\in\{1,\ldots,M_s\}$ は
長さ $D_s$ の対列 $\omega^m=(p_0^m,\ldots,p_{D_s-1}^m)$ であり，
公開 prefix に対して $p_j^m=p_j$（$j<\min(k,D_s)$）を要求する．
補完アルゴリズムは公開列と独立な探索 seed のみを引数に取る．

\textbf{実装記述．} 標準 \code{nextgen_safe_build} は
\begin{equation}\label{eq:profile}
 (D_s,w,M_s,Q,q_r)=(16,250,6,600000,1)
\end{equation}
であり，\code{legacy-fixed-six} を使う．基礎色順の添字 $0,1,2,3$ に対し
$(0,1,2,3)$，$(0,2,1,3)$，$(0,3,1,2)$，$(1,2,0,3)$，$(1,3,0,2)$，$(2,3,0,1)$
の各列を $2$ 組に分け，公開 prefix 後に反復する．設定 seed があるとシナリオ順と色対応を置換する．
別方式 \code{seeded-authoritative} は探索 seed と sample index から導いた seed を
\code{PuyoSequence} に渡す．名前の \code{authoritative} は乱数生成方式の再利用を指し，試合の実未来へのアクセスを意味しない．
特に $6$ 代表未来は真の条件付き分布からの独立標本ではない．

\subsection{ノード評価を与える方法}\label{sec:evaluator}
目標連鎖数 $L_*\ge1$ をあらかじめ固定する（標準 $10$）．
探索木のノード $v$ は，状態 $x(v)$，根行動 $a_0(v)$，配置列 $\mathbf a(v)$，
シナリオ番号，親ノード，直前遷移の $L,S,d$ を持つ有限レコードとする．
評価器 $U:\mathcal N\to\R$ の定義域 $\mathcal N$ はこのような有効ノード全体である．
ここでは評価器を探索アルゴリズムの固定パラメータとし，学習方策とは別に固定する．
現行版との対応を閉じるため，次の\textbf{実装による定義}を採用する．
$\psi(v)\in\R^{24}$ と $d_U(v)\in\{0,1\}$ を，固定 SHA の
\src{agents/chain_structure.py}{ChainStructureEvaluator} が，ノードの盤面・親の評価結果・直前行動を受け取り，
\code{_build_features}，\code{_action_features}，\code{_score} の順に計算する次表の値とする．
予算・重みは固定した設定ファイルの値を使う．$\psi$ は学習対象の未定義関数ではなく，この有限手続きの出力である．

\begin{longtable}{p{.26\linewidth}p{.67\linewidth}}
\toprule 評価特徴群（合計 24 成分） & 数学上の意味・算出対象\\\midrule\endhead
見込み連鎖数・得点 & 最大 3 個の仮想キーを補う有限探索で選ばれた発火の $L,S$．未発見は $0$．\\
必要キー数，発火点高さ，保護度 & 同じ仮想発火に必要な追加セル数，発火セル高さ，周囲の占有に基づく実装の保護度．未発見時は $0$．\\
残り 2 連結，残り 3 連結 & 仮想発火後の色成分のうち大きさが $2,3$ の個数．\\
連結辺 & 現盤面の同色色成分内の隣接辺数と，仮想発火後の同数の和．\\
接続候補，到達可能発火点 & \code{connection_candidates} の有限接続レコード数，到達可能な拡張セルを持つ大きさ $3$ の色成分数．\\
伸長箇所，土台セル，折返し空間 & 到達可能拡張セルの和集合の大きさ，$y\le2$ で下が支えられた色セル数，隣接列の可視空き高さの小さい方の総和（色セルなしなら $0$）．\\
凹凸，高さ幅，谷，突起 & 隣接列高さの絶対差の総和，最大と最小の高さ差，両隣より低い谷の深さ総和，内部列の両隣より高い突起の高さ総和．\\
危険度，おじゃま数，隠れ行占有数 & 式 \eqref{eq:danger} の値，盤面全体のおじゃま数，$y\ge12$ の占有数．\\
切断，浪費，発火点損傷，小発火 & 親に対する接続辺・接続候補の損失，隠れ行増加と未達発火の色資源損失，仮想発火能力の低下，$0<L<L_*$ の指示値．\\\bottomrule
\end{longtable}
構造評価用の色成分は，消去判定とは異なり $D$ 全域で隣接を定義する．
評価器内部は紫も含む 5 色の仮想列挙に対応するが，通常の公開盤面・ツモは本稿の 4 色である．
保護度・接続候補などの細かな列挙方法は上記関数をそのまま定義とする．
この表だけから独立実装を作り Python/native の全ビット一致を主張するものではない．
数学的モデルとしては，型を満たす別の明示した評価器にも置換できるが，その場合は別のモデル設定になる．

列 $x$ の最高占有セルの高さを $h_x=1+\max\{y:B(x,y)\ne\emptycell\}$ とし，空列は $h_x=0$ とする．
おじゃまセル数を $n_g$ とする．実装の危険度 $d_{\rm ratio}$ は
\begin{equation}\label{eq:danger}
 d_{\rm ratio}=\min\left(1,\ 0.55\frac{\max(h_2,h_3)}{12}
                       +0.35\frac{\max_xh_x}{12}+0.10\min(1,n_g/30)\right).
\end{equation}
評価器の death 判定には中央 \emph{2 列}の高さによる保守的条件も含まれる．
これをゲームの単一窒息セル判定と同一視しない．
仮想キー探索の内部上限は pattern $512$ ノード，resolution $96$ ノード，候補 $12$ 個である．
打切りで仮想発火が未発見なのか，この有限列挙を完了して未発見なのかも区別する．

\paragraph{評価特徴の数学的な展開．}
重み付けに用いる値を，さらに集合演算で記述する．
以下で $B$ はノード $v$ の盤面であり，構造用同色色成分全体を $\mathcal K_s(B)$ とする．
可視消去成分 $\mathcal K(B)$ と異なり，大きさ $1$ 以上，領域 $D$ 全体を対象とする．
占有セル集合を $O_B=\{u\in D:B(u)\ne\emptycell\}$ とする．
列番号を頂点，連続する列を辺とするグラフで，$h_x<12$ の列のみを残し，
列 $2$ または $3$ から到達する列の集合を $R_B$ とする．これは評価器の近似到達性であり，実時間 planner とは異なる．
成分 $K\in\mathcal K_s(B)$ の拡張セル集合 $X_B(K)$ と内部隣接辺数 $e(K)$ を
\begin{equation}\label{eq:extensions}
\begin{split}
 X_B(K)&=\{(x,h_x)\in D_v\setminus O_B:x\in R_B,\
                   \exists u\in K\text{ が }(x,h_x)\text{ に隣接}\},\\
 e(K)&=|\{\{u,v\}\subseteq K:u\text{ と }v\text{ は隣接}\}|
\end{split}
\end{equation}
とする．隣接は同じ上下左右規則を $D$ へ拡張したものである．
通常色セル数を $n_c(B)$，隠れ行占有数を $n_h(B)$，成分内辺総数を $e_B$ とする．
$\mathrm{color}(K)$ は成分 $K$ の共通色である．接続候補集合 $A_B$ は，
$u=(x,h_x)\in D_v\setminus O_B$，$x\in R_B$ に対し，同じ色 $c$ の異なる成分が $2$ 個以上隣接する
対 $(u,c)$ 全体とする．伸長数，発火可能な 3 連結数，土台セル数はそれぞれ
\begin{equation}\label{eq:structural-counts}
\begin{gathered}
 \left|\bigcup_{K\in\mathcal K_s(B)}X_B(K)\right|,\qquad
 \sum_{K\in\mathcal K_s(B)}\ind\{|K|=3,\ X_B(K)\ne\varnothing\},\\
 \sum_{\substack{(x,y)\in D\\B(x,y)\in\mathcal C}}
 \ind\{y\le2,\ y=0\text{ または }(x,y-1)\in O_B\}.
\end{gathered}
\end{equation}
となる．ここで $e_B=\sum_{K\in\mathcal K_s(B)}e(K)$，
$n_c(B)=\sum_{u\in D}\ind\{B(u)\in\mathcal C\}$，
$n_h(B)=\sum_{(x,y)\in D}\ind\{y\ge12,(x,y)\in O_B\}$ である．

可視余裕を $a_x=[12-h_x]_+$，隣列の高さを
$l_x=h_{x-1}$（$x>0$），$l_0=h_1$，$r_x=h_{x+1}$（$x<5$），$r_5=h_4$ とする．
折返し余裕 $f_B$，凹凸 $r_B$，高さ幅 $s_B$，谷 $w_B$，突起 $b_B$ は
\begin{equation}\label{eq:shape-features}
\begin{aligned}
 f_B&=\ind\{n_c(B)>0\}\sum_{x=0}^4\min(a_x,a_{x+1}),
 &r_B&=\sum_{x=0}^4|h_x-h_{x+1}|,\\
 s_B&=\max_xh_x-\min_xh_x,
 &w_B&=\sum_{x=0}^5[\min(l_x,r_x)-h_x]_+,\\
 b_B&=\sum_{x=1}^4[h_x-\max(l_x,r_x)]_+.
\end{aligned}
\end{equation}

仮想キー候補は $(A,c)$ とし，$A\subset D_v\setminus O_B$ は $1$ から $3$ 個の追加セル，
$c$ は全追加セルに共通の色である．各列の追加セルは $(x,h_x),(x,h_x+1),\ldots$ の連続 prefix とし，
$x\in R_B$ かつ高さ $12$ 未満を要求する．$B$ の $A$ へ色 $c$ を入れた盤面を $B\oplus(A,c)$ と書く．
追加セルに触れる大きさ $4$ 以上の同色可視成分があり，その中に元の色セルが少なくとも $1$ 個あることを要求する．
これを満たす最初の成分の元セル集合を anchor 集合 $A_0$ とする．
列順・高さ順の最初の追加セルから調べて同点を解消する．
どれか $1$ 列の最上追加セルを取り除いても同条件が成立する候補は除外する．
これが実装の「小さいキーを優先する」検査であり，任意の構成に対する大域的最小性の証明ではない．

有限予算内で解決した候補集合を $\mathcal Q_B$ とする．その列挙順は，追加個数昇順，
重複を許す列番号の昇順組合せを鏡映同値類にまとめた順，通常色順，同値類内の列列の辞書順である．
鏡映は $x\mapsto5-x$．mirror 対を予算途中で片方だけ追加せず，pattern/resolution の残予算を検査してから処理する．
実装は紫も試すが，4 色盤面に紫の既存 anchor はなく，3 個以下の追加だけでは上の条件を満たさない．
各候補について $\mathcal R(B\oplus(A,c))$ と式 \eqref{eq:score} の素点を計算する．
仮想発火には全消し保留ボーナスを加えない．

保護度 $p(A_0,A)$ は，anchor から可視領域内の隣接セルへ向かう向き付き辺の集合を $E_0$ として
\begin{equation}\label{eq:protection}
 p(A_0,A)=\begin{cases}
 \displaystyle\frac{|\{(u,v)\in E_0:v\in O_B\setminus A\}|}{|E_0|}&|E_0|>0,\\
 0&|E_0|=0
 \end{cases}
\end{equation}
である．候補の順位は，連鎖数降順，キー数昇順，素点降順，残り 3 連結数，2 連結数，内部辺数，伸長数，
保護度の各降順，発火高さ昇順，最後に canonical signature の SHA-256 の辞書式降順で決める．
canonical signature は色・鏡映標準化した追加セルと anchor・連鎖数・素点・キー数・発火高さ・残り連結と伸長数のレコードである．
同一 signature の候補は順位の高い方を保持し，完全同順位は列挙が早い方を採る．
$\mathcal Q_B$ の先頭を $q_B^*$ とする（空なら $\none$）．
見込み連鎖数 $\ell_B$，必要キー数 $k_B$ 等は $q_B^*$ の各値で，未発見は $0$ として重み付けする．
$\ell_B$ は実際の配置列で実現する連鎖数と区別する．

親盤面を $B_-$，直前配置の連鎖数を $L$，消去した色セル数を $n_v$ とする．
未達発火 bit を $p_L=\ind\{0<L<L_*\}$，目標達成 bit を $a_L=\ind\{L\ge L_*\}$ とする．
切断値 $t_B$ と浪費値 $u_B$ は
\begin{equation}\label{eq:action-costs}
\begin{split}
 t_B&=(1-a_L)\{[e_{B_-}-e_B]_+ + [|A_{B_-}|-|A_B|]_+\},\\
 u_B&=[n_h(B)-n_h(B_-)]_+
       +p_L\max([n_c(B_-)+2-n_c(B)]_+,n_v).
\end{split}
\end{equation}
発火点損傷は，親の仮想発火が存在し目標未達の場合にだけ計算する．
子の仮想探索が完了して発火未発見なら $\max(1,\ell_{B_-})$，
子にも仮想発火があれば $[\ell_{B_-}-\ell_B]_+ + [k_B-k_{B_-}]_+$，
それ以外は $0$ とする．親または行動の評価がない根ノードでは，これらの行動特徴を $0$ とする．


重みを $\mathbf w=(w_1,\ldots,w_{24})\in\R^{24}$，致命評価値を $u_{\rm fatal}=-10^{12}$ とする．
$d_U(v)=1$ は実装の death，探索完了時の発火点未発見，構造的行止まり，直前行動の death のいずれかが成立する場合である．
評価値を
\begin{equation}\label{eq:evaluator}
 U(v)=\begin{cases}u_{\rm fatal}&d_U(v)=1,\\
 \displaystyle\sum_{j=1}^{24}w_j\psi_j(v)&d_U(v)=0
 \end{cases}
\end{equation}
とする．$\psi_j$ は表の順に展開した第 $j$ 成分である．標準重みの正本は
\src{train/config/nextgen_chain_structure.yaml}{nextgen_chain_structure.yaml} であり，
例えば見込み連鎖数 $45000$，見込み得点 $0.25$，危険度 $-480000$，小発火 $-16000$ である．
仮想キーの $L$ は実際の公開ツモで達成した連鎖数ではない．
得点目標モードは別の \code{nextgen_score_structure.yaml} を使う．

\subsection{幅制限探索と根の保護}\label{sec:beam}
1 つのシナリオ $m$ を固定する．深さ $d$ の生存ノード列を $\mathcal L_d^m$，
次層の展開候補列を $\mathcal E_{d+1}^m$ とする．$\mathcal L_0^m$ は初期状態だけからなる．
既定では連鎖数 $L\ge1$ の子を「発火終端」として記録して以後の展開から外し，
連鎖なしで窒息する子も外す．「発火終端」は探索の停止であって試合終了ではない．
展開は
\begin{equation}\label{eq:beam-expand}
 \mathcal E_{d+1}^m=\bigl(T(x(v),p_d^m,a):
       v\in\mathcal L_d^m,\ a\in\mathcal A_{\rm geo}(B(v))\bigr)
\end{equation}
から有効な子ノードを構成し，選択中テンプレートの構造制約を検査する操作である．
$B(v)$ は $x(v)$ の盤面成分である．括弧は順序付き列挙を表し，展開予算 $Q$ に達したら途中で止める．
根では前節の到達可能性も最終候補化時に検査する．

生存候補を $U$ の降順で並べ，同点では根行動 ID，状態の標準 byte 列，直前行動 ID，経路 ID 列の昇順で並べる．
この全順序を $\prec_U$ とする．状態が合流した場合も根行動が異なる候補は分ける．
同じ根と同じ compact 遷移キー（盤面・全消し等の状態，scenario，ツモ cursor，深さ）の候補が複数あれば，
$\prec_U$ の先頭だけを保持する．これを重複除去と定義する．

切詰め写像 $\operatorname{Keep}_{w,q_r}$ は，根行動 ID 順に各根の第 $1$ 候補を取り，
次に各根の第 $2$ 候補という順で $q_r$ 巡まで確保し，残りを $\prec_U$ 順で埋める．
どの段階でも総数 $w$ を超えない．該当候補がなければ飛ばす．従って
\begin{equation}\label{eq:beam-prune}
 \mathcal L_{d+1}^m=\operatorname{Keep}_{w,q_r}
     (\text{重複除去後の非発火・非窒息候補}),\qquad |\mathcal L_{d+1}^m|\le w.
\end{equation}
単純に全体の上位 $w$ 個を選ぶ式とは異なる．$w$ が小さければ全根の確保も保証しない．
各シナリオを固定順で処理し，全シナリオで展開予算 $Q$ を共有する．
盤面評価器の仮想キー探索予算はこれとは別の内部上限である．

\subsection{発火分類とシナリオ集約}\label{sec:aggregate}
既定した目標連鎖数 $L_*$ に加え，任意の得点目標を $S_*\in\{1,2,\ldots\}\cup\{\none\}$ とする．
得点のみを目標にする flag を $b_S\in\{0,1\}$，強制安全対応の文脈 flag を $b_F\in\{0,1\}$ とする．
$b_S=1$ なら $S_*\ne\none$ を要求する．発火分類 $c(L,S)$ は次の先勝ち条件で決まる．
\begin{equation}\label{eq:fireclass}
 c(L,S)=\begin{cases}
 \mathsf{quiet}&L=0,\\
 \mathsf{winning}&S_*\ne\none\text{ かつ }S\ge S_*,\\
 \mathsf{target}&b_S=0\text{ かつ }L\ge L_*,\\
 \mathsf{forced}&b_F=1,\\
 \mathsf{premature}&\text{それ以外}.
 \end{cases}
\end{equation}
$\mathsf{winning}$ は閾値到達の内部名であり，対戦勝利の証明ではない．
未評価を $\mathsf{unavailable}$ とし，優先度写像 $P$ は
$\mathsf{unavailable},\mathsf{premature},\mathsf{quiet},\mathsf{forced},\mathsf{target},\mathsf{winning}$
の順に $0,1,2,3,4,5$ を割り当てる．

発火子 $v$ の $L,S$ に対し，不足量 $g(v)$ を $b_S=0$ なら $[L_*-L]_+$，
$b_S=1$ なら $[S_*-S]_+/1000$ とする．不足罰則係数を $\lambda_g\ge0$（標準 $1000$），
許容発火分類集合を $\mathcal C_+=\{\mathsf{forced},\mathsf{target},\mathsf{winning}\}$ とする．
発火評価値 $U_{\rm fire}$ は
\begin{equation}\label{eq:firevalue}
 U_{\rm fire}(v)=U(v)-\lambda_g g(v)\ind\{c(L,S)=\mathsf{premature}\}
                     +S\ind\{c(L,S)\in\mathcal C_+\}.
\end{equation}

各根 $a$ の各シナリオ $m$ について，実際に根を評価した bit を $e_{am}$，
選ばれた代表分類を $c_{am}$，発火代表の得点・連鎖数を $s_{am},\ell_{am}$，
発火評価を $u_{am}$，非発火継続評価を $v_{am}$ とする．
数値が存在しなければその値は $\none$ とする．代表は
\src{agents/long_horizon_search.py}{_ScenarioTracker} の固定順位で選んだ観測済み候補である．
これはシナリオ内の全合法経路の最大値ではない．
評価済み添字集合 $I_a=\{m:e_{am}=1\}$，coverage $\rho_a=|I_a|/M_s$ とする．
$I_a$ が空なら分類は $\mathsf{unavailable}$．それ以外では
$c_a$ を $\{c_{am}:m\in I_a\}$ のうち $P$ が最大の分類とし，その支持数を
$n_a=|\{m\in I_a:c_{am}=c_a\}|$ とする．

$m\in I_a$ に対し $\tilde s_{am},\tilde\ell_{am}$ は，$c_{am}=c_a$ で発火代表がある場合に
$s_{am},\ell_{am}$，それ以外では $0$ とする．有限数列 $z=(z_1,\ldots,z_n)$ の
平均と標準偏差を $\mu(z)=n^{-1}\sum_jz_j$，
$\sigma(z)=\sqrt{n^{-1}\sum_j(z_j-\mu(z))^2}$ と定義し，空列では両者 $0$ とする．
$c_a=\mathsf{quiet}$ なら候補値 $V_a$ は評価済み非発火継続の $v_{am}$ の平均，
発火分類なら同じ分類の発火代表が持つ $u_{am}$ の和とする．該当数値がなければ $V_a=-10^{12}$ とする．
また $k_a$ を発火が一度でも観測されたシナリオ数，$j_a$ を $\mathsf{quiet}$ のシナリオ数，
$\bar v_a$ を有効な非発火継続評価の平均（なければ $-\infty$）とする．
$-\infty$ は全実数より小さい順序用の値であり，学習特徴には渡さない．

根の基本順位キー $K_a$ は，次の列の辞書式降順である．$\mathrm{id}(a)$ は配置 ID とする．
\begin{equation}\label{eq:rootrank}
\begin{split}
 K_a=(&P(c_a),\rho_a,n_a,V_a,
 \textstyle\sum_{m\in I_a}\tilde s_{am},\sum_{m\in I_a}\tilde\ell_{am},k_a,\\
 &\min_{m\in I_a}\tilde s_{am},\min_{m\in I_a}\tilde\ell_{am},
 -\sigma(\tilde s_a),-\sigma(\tilde\ell_a),\bar v_a,j_a,-\mathrm{id}(a)).
\end{split}
\end{equation}
$\tilde s_a,\tilde\ell_a$ は $I_a$ をシナリオ ID 順に並べた列で，この式の空集合の最小値は $0$ とする．
実装では浮動小数点の左からの加算順も固定する．

標準 $6$ シナリオで全根が評価済み，安全構築文脈，発火で展開停止，
winning/forced の根なしという条件の下で，目標到達の支持が $2$ 未満の根と，
全 $6$ シナリオで打切りなく深さ $D_s$ に達した非発火生存根が共存するときは補正する．
後者は深さ $D_s$ に残存候補があり，各継続評価が有限で $u_{\rm fatal}$ より大きいことも要求する．
このとき前者の $P(c_a)$ を $2.25$，後者を $2.5$ に置換する．
これが PR \#188 に含まれる \code{weak_scenario_support} 補正である．
$\rho_a<1$ は未評価を含むためであって失敗確率ではない．
$V_a$ や平均連鎖数は有限仮説に対する探索統計であり，真の勝率や不偏期待値ではない．

\section{テンプレート，戦術，実行の分離}\label{sec:tactics}
\subsection{テンプレートの制約}
テンプレート記号の有限集合を $\Sigma$，底面を基準とした対象セルから記号への部分写像を
$P_T:D\rightharpoonup\Sigma$ とする．$\rightharpoonup$ は一部の入力でのみ定義される写像を意味する．
色束縛 $\beta:\Sigma\to\mathcal C$ は，同色関係 $R_=\subseteq\Sigma^2$ と異色関係
$R_{\ne}\subseteq\Sigma^2$ に対して，それぞれ $\beta(A)=\beta(B)$ と $\beta(A)\ne\beta(B)$ を満たす．
異なる記号であるだけでは異色を要求しない．
空必須セル集合 $D_0$，占有必須セル集合 $D_+$，禁止色条件集合 $D_-\subseteq D\times\Sigma$ も与える．
これらをまとめたレコードをテンプレート $T$ と呼び，$\operatorname{dom}P_T$ を $P_T$ の定義域とする．
完全成立条件 $\operatorname{Sat}(B,T,\beta)$ は
\begin{equation}\label{eq:template}
 \begin{gathered}
 B(u)=\beta(P_T(u))\ (u\in\operatorname{dom}P_T),\qquad
 B(u)=\emptycell\ (u\in D_0),\qquad B(u)\ne\emptycell\ (u\in D_+),\\
 B(u)\ne\beta(A)\ ((u,A)\in D_-)
 \end{gathered}
\end{equation}
を全て満たすこととする．
\code{.} と \code{*} は $\operatorname{dom}P_T$ に含まれない判定対象外のセルである．
鏡映・平行移動・底上げは，これらのセル集合を先に変換してから同じ述語で評価する．

variant の重みを $w_T>0$，同じ template 内の全 variant の重み和を $W_T$，
記号・空・占有の必須条件数を $n_T>0$ とする．部分盤面で充足した必須条件数を $n_{\rm ok}$，
既知セルによる矛盾と禁止色違反の個数を $n_{\rm conflict}$ とする．未観測セルは充足に数えない．
静的適合値は
\begin{equation}\label{eq:static-fit}
 f_{\rm static}=\frac{w_T}{W_T}\frac{n_{\rm ok}-n_{\rm conflict}}{n_T}
\end{equation}
である．現行 matcher の最終 score は，この静的評価と公開ツモ prefix 上の有限探索結果から
\src{agents/template_catalog.py}{match_templates} が計算する値と定義する．
静的 score と探索後の score，部分適合 \code{fit} と完全成立 \code{complete} は異なる．
公開ツモだけを用いた具体的配置列を適合の根拠とし，sampled future を確定根拠としない．

template ごとの最良 variant の score を $f_j$，候補番号を $j\in J_T$ とする．
$J_T$ は選択対象 template の有限添字集合である．初期選択では enabled 全候補，
対応・発火後の再選択では fit 根拠がある候補だけを用いる．
選択は score 最大（固定 key で同点解消），または温度 $\tau_T>0$ による
\begin{equation}\label{eq:template-softmax}
 p_T(j)=\frac{\exp((f_j-f_{\max})/\tau_T)}{\sum_{k\in J_T}\exp((f_k-f_{\max})/\tau_T)},
 \qquad f_{\max}=\max_{k\in J_T}f_k
\end{equation}
とする．再選択で $J_T=\varnothing$ なら自由構築へ移る．この初期形状選択は戦術 RL の出力ではない．

テンプレート phase を $\varphi=(j,\beta,n,N,\mathrm{active})$ とする．
$j$ は選択済み template/variant/変換，$n$ は採用されたテンプレート構築 decision 数，
$N\ge1$ は上限（標準 $14$），$\mathrm{active}$ は有効 bit である．
同じ組ぷよの再要求や採用失敗では $n$ を増やさない．採用した \code{build_template} で $1$ 増やし，
$n=N$，完成，整合喪失，他戦術への切替で閉じる．相殺・カウンター・発火の完了は実行 event で確認してから再選択する．

\subsection{戦術と構造的な選択可能性}
戦術集合 $\mathcal Z$ を，次表の順序を持つ $6$ 個の識別子集合とする．
\begin{center}\begin{tabular}{ll}\toprule ID & 目的\\\midrule
\code{build_main}&大連鎖本線の構築\\
\code{build_template}&選択済み定型の構築\\
\code{fire_main}&本線の発火準備・発火\\
\code{cancel}&到来する攻撃の相殺準備・相殺\\
\code{counter}&最初のおじゃま落下後の次の 1 配置で発火\\
\code{decisive_short_attack}&決定的機会を狙う短期攻撃\\\bottomrule\end{tabular}\end{center}

候補 $c$ は，根行動 $a(c)\in\mathcal A$，組ぷよと配置の有限計画列，
所属戦術集合，予測の前提，根の合法性・到達可能性，評価根拠を持つレコードである．
候補集合を $\mathcal C_n$，戦術 $z$ の候補を探索終了時の固定順で並べた列を $\mathcal C_n(z)$ とする．
添字 $n$ は同一プレイヤーの意思決定番号であり，tick $t$ とは異なる．

各数値根拠を $(v,e,s)$ と定める．$v\in\R\cup\{\none\}$ は値，
$e$ は evaluated，partial，not\_evaluated，unsupported のいずれか，
$s$ は visible\_exact，public\_estimate，sampled\_future のいずれかである．
未評価・未対応では $v=\none$ とする．partial の有限値は評価済み範囲での値であり確定保証ではない．
生存率推定，期限，火力にもこの型を用いる．

全候補の根が合法・到達可能であり，戦術固有の構造条件を満たすよう provider が列を作る．
構造条件には有効 template phase，対象攻撃，期限内到達，落下後の発火点が含まれる．
火力不足や「勝ちそうにない」という選好だけで候補を隠さない．戦術 mask $m_n(z)$ と代表候補 $c_n^*(z)$ は
\begin{equation}\label{eq:tactic-mask}
 m_n(z)=\ind\{|\mathcal C_n(z)|>0\},\qquad
 c_n^*(z)=\text{$\mathcal C_n(z)$ の先頭}\quad(m_n(z)=1)
\end{equation}
とする．探索打切りによる未発見は \code{not_found_within_budget} と記録する．
$c_n^*$ は選択器から変更できず，戦術選択の後に再探索もしない．

期限予測は点ではなく区間を持てる．発火開始区間を $[t_{s,-},t_{s,+}]$，完了区間を
$[t_{e,-},t_{e,+}]$，受け手の落下境界区間を $[t_{d,-},t_{d,+}]$ とする．
これらの端点は公開推定または公開事実として出典を持つ非負整数である．
$t_{e,+}\le t_{d,-}$ は保守的な相殺間に合い条件だが，必要条件ではない．
同じ配置境界で攻撃生成と落下が処理される場合，時刻推定の誤差は相関するためである．
現行 selector は，公開根拠を持つ 1 手完全相殺で開始期限を満たし，完了と落下の区間端点が一致する場合も採用できる．
カウンターの根拠は「落下後いつか」ではなく，最初の落下後の最初の配置解決で $L>0$ となる計画である．

\subsection{現在のルール選択器}
優先順列を
$\mathbf p=(\mathsf{cancel},\mathsf{counter},\mathsf{short},\mathsf{template},\mathsf{fire},\mathsf{build})$
とする．ここでは表の長い ID を順に略記した．
$e_n(z)\in\{0,1\}$ を \code{RuleTacticSelector.select} の適格判定とする．
その条件は，相殺可能な即時脅威または前述の 1 手完全相殺，カウンターは即時脅威，
短期攻撃は送信量 $30$ 以上かつ相手占有セル $48$ 以上，テンプレートは active phase，
本線発火は予測 $10$ 連鎖以上（得点目標指定時はその閾値以上）かつ fatal rate $0$，自由構築は常に真である．
即時脅威で完全相殺不能なら，相殺・カウンターの候補を fatal rate 昇順，相殺量・送信量降順，
sampled な生存根拠への罰，元優先順で比較して適格側を選び直す．欠損 fatal rate は $1$ と扱う．
この手順で定義された $e_n$ を使い，現行戦術 $z_n^{\rm rule}$ は
\begin{equation}\label{eq:rule}
 z_n^{\rm rule}=\text{$\mathbf p$ の順で最初の $z$：$m_n(z)e_n(z)=1$}
\end{equation}
とする．該当要素なしなら scheduler に処理を戻す．
この $e_n$ と構造 mask $m_n$ を混同しない．将来の学習器へルール優先順位を hard mask として持ち込まない．

\subsection{実行境界}
要求 tick を $t_n^{\rm req}$，完了通知受領 tick を $t_n^{\rm done}$，設定遅延を $\delta_n\in\N$，
timeout tick を $t_n^{\rm timeout}$ とする．予定採用 tick $t_n^{\rm act}$ は
\begin{equation}\label{eq:activation}
 t_n^{\rm act}=\begin{cases}
 t_n^{\rm req}+\delta_n&\text{configured mode},\\
 \max(t_n^{\rm done},t_n^{\rm req}+\delta_n)&\text{measured mode}.
 \end{cases}
\end{equation}
$t_n^{\rm act}>t_n^{\rm timeout}$，古くなった盤面，同一性不一致では要求を採用しない．
要求行動 $a_n^{\rm req}=a(c_n^*(z_n))$ と，実際の行動 $a_n^{\rm exec}$ は別々に記録する．
複数手の計画は根拠であり，実行するのは根の $1$ 手だけである．次の公開状態で再計算する．
合法根があるのに全戦術 mask が偽になるのは契約違反であり，通常は
\code{build_main} に決定的候補を用意する．操作不能・終端では方策を呼び出さない．
候補自体の fallback 印と，scheduler が提案を採用せず代替行動を実行する fallback outcome は別概念である．

\section{戦術選択を学習するモデル}\label{sec:learning}
本節の方策ネットワーク，損失関数，更新則は\textbf{提案}である．
公開契約・特徴量・ログを利用するが，PR \#188 に同じ学習器が存在すると主張しない．

\subsection{方策に渡す特徴量}
意思決定 $n$ における公開履歴と phase，候補 batch から生成した raw 要約を
$y_n=(y_{n1},\ldots,y_{nF})\in(\R\cup\{\none\})^F$ とする．
$F=93$ は \code{FEATURE_REGISTRY} の登録数である．
要約生成写像は，固定 SHA の \code{feature_summaries} と \code{TimingSummary.feature_summaries} が
登録名へ値を対応させる写像として定義する．入力にない・未評価の登録項目は $\none$ とする．
生盤面や raw tick を 93 次元の中へ追加しない．登録順と意味は次の通りである．
\begin{center}\begin{tabular}{p{.59\linewidth}rr}\toprule 要約群 & 個数 & 正規化 scale\\\midrule
脅威 none/potential/pressing/immediate の各指示値 &4&1\\
対応 possible/marginal/impossible/unknown の各指示値 &4&1\\
自分の発火中に相手が完了する配置数の下限・上限 &2&3\\
自他 $\times$ 本線・副砲 $\times$ 連鎖数・火力・危険度・発火点生存 &16&19/360/1/1\\
phase active・残数比・進捗・fit/no\_fit/unknown &6&1\\
前回戦術の各指示値 &6&1\\
同じ戦術の継続数 &1&14\\
各戦術 $\times$ available・known\_witness・候補数・火力・期限差の負/零/正・coverage・fatal rate &54&注記\\\bottomrule\end{tabular}\end{center}
最後の 9 項目の scale は順に $1,1,8,360,1,1,1,1,1$ である．
火力だけは対数変換，ほかは線形変換を用いる．
火力はおじゃま個数，危険度は $[0,1]$，発火点生存は指示値である．
期限差は公開推定の落下境界下限から発火完了上限を引いた値の符号，
coverage は評価済みシナリオ数の要求数に対する比，fatal rate は評価対象で観測した失敗の比である．
各比には評価範囲と根拠を sidecar に保持する．

脅威階級は，公開推定の落下境界までに残る操作回数が $1$ 以下なら immediate，$2$ 以下なら pressing，
それ以外なら potential とする．未評価時は none へ丸めず欠損にする．
対応階級と先打ち損失の区間は，公開応答計画の発火・操作再開の区間を
\src{puyo_env/nextgen_public_snapshot.py}{derive_timing_summary} で比較した結果として定義する．
情報がなければ対応は unknown，損失は欠損になる．
登録された自他本線・副砲の全項目が現行 producer で必ず埋まるわけではない．

特徴 $j\in\{1,\ldots,F\}$ の scale を $s_j^{\rm feat}>0$，変換種別を $\tau_j\in\{\mathsf{linear},\mathsf{log}\}$ とする．
欠損 bit $b_{nj}^{\rm miss}$，正規化値 $v_{nj}$，actor 入力 $x_n^{\rm pol}$ を
\begin{align}
 b_{nj}^{\rm miss}&=\ind\{y_{nj}=\none\},\label{eq:missing}\\
 v_{nj}&=\begin{cases}
 0&y_{nj}=\none,\\
 \operatorname{clip}([y_{nj}]_+/s_j^{\rm feat};0,1)&\tau_j=\mathsf{linear},\ y_{nj}\ne\none,\\
 \operatorname{clip}\left(\dfrac{\log(1+[y_{nj}]_+)}{\log(1+s_j^{\rm feat})};0,1\right)&\tau_j=\mathsf{log},\ y_{nj}\ne\none,
 \end{cases}\label{eq:normalize}\\
 x_n^{\rm pol}&=(v_{n1},\ldots,v_{nF},b_{n1}^{\rm miss},\ldots,b_{nF}^{\rm miss})\in[0,1]^{186}
 \label{eq:actor-input}
\end{align}
とする．mask は別に $m_n\in\{0,1\}^6$ として渡す．実際の学習入力はこの組だけとし，
非公開 seed，配牌 ID，候補列，生盤面，oracle 評価を追加しない．

\subsection{mask 付き方策と配置方策への合成}
学習パラメータ空間を $\Theta=\R^p$（$p\ge1$ はパラメータ数）とし，$\theta\in\Theta$ を actor のパラメータとする．
$f_\theta:[0,1]^{186}\to\R^6$ は戦術ごとの実数 logit（正規化前の score）を返すネットワークである．
例えば隠れ層幅を正整数列 $(d_1,\ldots,d_K)$，$d_0=186$，$d_{K+1}=6$ と固定し，
行列 $W_k\in\R^{d_k\times d_{k-1}}$ とバイアス $b_k\in\R^{d_k}$ をパラメータに含める．
$h_0=x$，$h_k=\max(0,W_kh_{k-1}+b_k)$（成分ごとの最大値，$1\le k\le K$），
$f_\theta(x)=W_{K+1}h_K+b_{K+1}$ とすれば具体的な多層パーセプトロンになる．
隠れ幅は実験設定として保存し，本稿では未実測の特定値を既定採用しない．

選択可能戦術集合を $\mathcal Z_n=\{z\in\mathcal Z:m_n(z)=1\}$ とする．
$\mathcal Z_n\ne\varnothing$ のとき，戦術方策 $\pi_\theta$ を
\begin{equation}\label{eq:policy}
 \pi_\theta(z\mid x_n^{\rm pol},m_n)=
 \begin{cases}
 \dfrac{\exp f_\theta(x_n^{\rm pol})_z}
       {\sum_{u\in\mathcal Z_n}\exp f_\theta(x_n^{\rm pol})_u}&z\in\mathcal Z_n,\\
 0&z\notin\mathcal Z_n
 \end{cases}
\end{equation}
で定義する．添字 $z$ は登録された戦術順での成分を意味する．数値計算では有効 logit の最大値を分子・分母の指数から引く．
全 false なら softmax を計算せず，前節の scheduler 境界で処理する．

探索 batch を固定したとき，要求配置の条件付き分布 $\Pi_\theta$ は
\begin{equation}\label{eq:composed-policy}
 \Pi_\theta(a\mid x_n^{\rm pol},m_n,\mathcal C_n)
 =\sum_{z\in\mathcal Z_n}\pi_\theta(z\mid x_n^{\rm pol},m_n)
       \ind\{a(c_n^*(z))=a\}.
\end{equation}
同じ配置に複数戦術が対応してよい．戦術 ID を抽選・記録して更新するため，
この配置確率を戦術の log-probability の代わりに使わない．
phase と前回戦術・継続数が入力にあるため，戦術の維持・切替を学習できる．
学習器は具体配置，目標連鎖数，探索幅，候補の順位，テンプレートの初期選択を直接出力しない．
各 decision で再選択する構成なので，別の学習可能な option 終了確率を暗黙に追加しない．

\subsection{履歴を状態とする可変時間モデル}\label{sec:semimarkov}
相手の方策 $\pi^-$，環境初期分布 $\mu_0$，探索・scheduler・規則の設定を 1 回の rollout 中は固定する．
相手の内部記憶と両者の scheduler 状態も含む拡張全状態を $\tilde s$ とし，その集合を $\widetilde{\mathcal S}$ とする．
乱数 stream を初期状態に含めれば進行は決定的であり，$\mu_0$ と戦術抽選によって軌跡の確率法則が定まる．

同一プレイヤーの意思決定 epoch を tick 列 $t_0<t_1<\cdots$ とする．
再試行が同一 tick に複数ある場合は event としてまとめ，学習上の時間を偽って増やさない．
$t_n$ から $t_{n+1}$ までの秒数を $\Delta_n=(t_{n+1}-t_n)/60>0$ とする．
意思決定間に相手や連鎖・おじゃま・fallback は進行する．最後は次 decision の代わりに終端または収集境界 tick を使う．
拡張状態から，選択戦術 $z$ の候補を実行して次 epoch に至る条件付き分布を
\begin{equation}\label{eq:kernel}
 K_\eta^{\pi^-}(\tilde s',\Delta\mid\tilde s,z)
 =\Pr\{\tilde S_{n+1}=\tilde s',\ \Delta_n=\Delta\mid\tilde S_n=\tilde s,Z_n=z\}
\end{equation}
と定義する．大文字 $\tilde S_n,Z_n$ は軌跡上の確率変数である．
この核は $F_\eta$，固定探索器，実行 planner，scheduler，相手方策の合成から誘導され，
独立に「学習済み環境モデル」を仮定しない．遷移に可変時間が付く制御モデルを半マルコフモデルという．
プレイヤーが全状態を観測できないため，ここでは部分観測の半マルコフ意思決定問題となる．

epoch $n$ までの公開履歴と自分の制御記録を $H_n$ とする．この拡張した履歴型も $\mathcal H$ に含める．
信念状態 $b_n(\tilde s)=\Pr(\tilde S_n=\tilde s\mid H_n)$ は理論上の条件付き分布であり，
現行実装がこれを推定しているとは主張しない．
$H_n$ を状態とすれば過去の全情報を保持するが，$x_n^{\rm pol}$ はその圧縮である．
従って $x_n^{\rm pol}$ だけでマルコフ性が成立する保証はない．後述の価値ネットワークは履歴価値の近似である．

\subsection{報酬・割引・価値の定義}
1 秒あたりの割引率を $\gamma\in(0,1]$ とし，decision 間の割引を $\Gamma_n=\gamma^{\Delta_n}$ とする．
対戦終了結果以外に基礎報酬を与えない．epoch 区間内で tick $t_T$ に終端へ到達した場合，
プレイヤー $i$ の区間基礎報酬を
\begin{equation}\label{eq:base-reward}
 R_n^0=\gamma^{(t_T-t_n)/60}z_i(s_{t_T})
\end{equation}
とし，終端がない区間は $R_n^0=0$ とする．この定義では「終端の $\pm1$」を実際の発生時刻で割り引く．
区間末尾報酬を割り引かず $\pm1$ とする別慣習と混ぜない．$\gamma<1$ は勝敗に時間選好も加えるため，
純粋な勝率最大化と同一ではない．$\gamma=1$ の有限対局で初めて期待勝敗との一致が成立する．

補助 potential を，有界写像 $\Phi:\mathcal H\to\R$ として事前に固定し，真の終端では $\Phi=0$ とする．
実装案では公開特徴だけの有界関数を用いる．係数 $\alpha\ge0$ を固定し，区間報酬を
\begin{equation}\label{eq:potential}
 R_n=R_n^0+\alpha\{\Gamma_n\Phi(H_{n+1})-\Phi(H_n)\}
\end{equation}
とする．初期検証は $\alpha=0$．後で補助を導入する場合は関数・係数を manifest に保存する．
小連鎖回数や全消し達成だけへの固定加点はこの目的に含めない．

以下の有限対局の目的では，対象の方策・初期分布で対局が確率 $1$ で有限時間に終わると仮定する．
実装で時間制限に達したデータを，この仮定が検証できた証拠にはしない．
終端までの区間数を $N$，epoch $n$ からの収益を $G_n$ とする．
\begin{equation}\label{eq:return}
 G_n=\sum_{k=n}^{N-1}\left(\prod_{j=n}^{k-1}\Gamma_j\right)R_k,
 \qquad G_N=0.
\end{equation}
学習目的は $J(\theta)=\mathbb E_{\mu_0,\pi^-,\pi_\theta}[G_0]$ の最大化である．
期待値には組ぷよ・おじゃま列抽選・相手行動・探索乱数・戦術抽選を含める．
任意の固定戦術方策を $\pi$ とし，その実験で正の確率を持つ履歴と選択戦術の対を $(H,z)$ とする．
履歴価値と行動価値は
\begin{equation}\label{eq:value}
\begin{aligned}
 V^\pi(H)&=\mathbb E[G_n\mid H_n=H],\\
 Q^\pi(H,z)&=\mathbb E[G_n\mid H_n=H,Z_n=z],\\
 A^\pi(H,z)&=Q^\pi(H,z)-V^\pi(H).
\end{aligned}
\end{equation}
である．
これらは履歴に対する定義であり，特徴要約だけに厳密な Bellman 方程式を仮定しない．

\begin{proposition}[potential の相殺]\label{prop:potential}
有限対局で式 \eqref{eq:return} の補助部分は $-\alpha\Phi(H_0)$ になる．
\end{proposition}
\begin{proof}
$D_0=1$，$D_k=\prod_{j=0}^{k-1}\Gamma_j$ と置くと，$D_{k+1}=D_k\Gamma_k$ である．
従って補助和は
$\alpha\sum_{k=0}^{N-1}\{D_{k+1}\Phi(H_{k+1})-D_k\Phi(H_k)\}
=\alpha\{D_N\Phi(H_N)-\Phi(H_0)\}=-\alpha\Phi(H_0)$．
\end{proof}
初期分布が方策に依存しなければ目的間の差は定数である．
これは無限時間・有限標本・関数近似下で学習が必ず安定するという主張ではない．
時間制限では終端 potential を強制的に $0$ にせず，続きの価値を bootstrap する．

\subsection{教師による初期化と PPO 更新}
教師データ集合を $\mathcal D_{\rm BC}=\{(x_j,m_j,z_j)\}_{j=1}^{M}$ とする．
$M\ge1$ はサンプル数，$x_j$ は保存済み 186 次元入力，$m_j$ は同時点の mask，
$z_j$ はその mask 内のルール戦術である．$j$ はデータ番号であり探索深さではない．
交差エントロピー損失
\begin{equation}\label{eq:bc}
 \mathcal L_{\rm BC}(\theta)=-\frac1M\sum_{j=1}^M\log\pi_\theta(z_j\mid x_j,m_j)
\end{equation}
を最小化して初期化する．候補再順位付けや配置パラメータの教師損失は持たせない．
episode と配色系列単位で train/validation/test を分離し，同じ軌跡の decision を別 split にしない．

価値近似器を $V_\nu:[0,1]^{186}\to\R$，そのパラメータを $\nu$ とする．
方策 \eqref{eq:policy} の出力層を $1$ 次元にした別ネットワークとして具体化できる．
rollout 収集時のパラメータを $\theta_{\rm old},\nu_{\rm old}$ とし，更新中に固定する．
真の終端 bit を $d_n^T$，時刻差付き trace 率を $\Lambda_n=\lambda^{\Delta_n}$ とする．
$\lambda\in(0,1]$ は 1 秒あたりの trace 減衰率である．
次価値 $\bar V_{n+1}$ は，真の終端なら $0$，それ以外では実際の境界特徴が取得できる場合の
$V_{\nu_{\rm old}}(x_{n+1}^{\rm pol})$ と定義する．境界特徴がなければその行は value 更新の対象外である．
rollout 末尾の次添字の advantage は $0$ と定義しておく．
TD 残差 $\delta_n^{\rm TD}$，GAE 近似 $\hat A_n$，価値教師 $\hat G_n$ を
\begin{align}
 \delta_n^{\rm TD}&=R_n+\Gamma_n\bar V_{n+1}
                            -V_{\nu_{\rm old}}(x_n^{\rm pol}),\label{eq:td}\\
 \hat A_n&=\delta_n^{\rm TD}+\Gamma_n\Lambda_n c_n^{\rm trace}\hat A_{n+1},
 \qquad \hat G_n=\hat A_n+V_{\nu_{\rm old}}(x_n^{\rm pol})\label{eq:gae}
\end{align}
とする．$c_n^{\rm trace}=1$ は同じ完全な軌跡の次区間をこの rollout 内で利用できる場合，
それ以外は $0$ である．真の終端では $d_n^T=1,c_n^{\rm trace}=0$ とし，次価値項自体を評価しない．
rollout 末尾・時間制限では $d_n^T=0,c_n^{\rm trace}=0$ とし，実際の境界観測から次価値を評価する．
$\hat A$ は末尾から逆順に計算する．終端後のダミー特徴を実測観測と誤認しない．
必要な次特徴が取得できなければ，その行を bootstrap 可能と偽らず value 不適格とする．
初期案 $\gamma=0.99,\lambda=0.95$ は 1 秒あたりの実験値であり，1 decision あたりの定数ではない．

actor 適格添字集合を $I_{\rm act}$，value 適格添字集合を $I_{\rm val}$ とする．
前者には，同じ batch・mask に対する行動が実際に activated され，保存 log-probability が検証できる行だけを含める．
後者には，時間差・報酬・次価値・継続履歴が再構成できる行を含める．両者は一致しなくてよい．
stale/timeout/scheduler fallback は actor から除くが，経過時間と報酬は軌跡から削除しない．
欠落した区間を飛び越えて $\Delta_n$ を縮めない．これは介入を含む実行システムを学習する仕様であり，
介入のない方策の収益に対する不偏推定を保証するものではない．

$n\in I_{\rm act}$ の選択戦術を $z_n$，保存された旧方策 log-probability を
$\ell_n^{\rm old}=\log\pi_{\theta_{\rm old}}(z_n\mid x_n^{\rm pol},m_n)$ とする．
以下の平均は適格集合が非空の場合に定義し，空の場合は対応する損失項を $0$ と定めて更新をスキップする．
確率比 $r_n(\theta)$ と clipping 幅 $\epsilon\in(0,1)$ を用い，PPO の方策目的を
\begin{align}
 r_n(\theta)&=\exp\{\log\pi_\theta(z_n\mid x_n^{\rm pol},m_n)-\ell_n^{\rm old}\},\label{eq:ratio}\\
 \mathcal L_{\rm clip}(\theta)&=\frac1{|I_{\rm act}|}\sum_{n\in I_{\rm act}}
 \min\{r_n(\theta)\hat A_n,
       \operatorname{clip}(r_n(\theta);1-\epsilon,1+\epsilon)\hat A_n\}
 \label{eq:ppo}
\end{align}
と定め，最大化する\cite{ppo}．同じ保存済み mask で新旧確率を計算する．
更新時に新たな探索で mask を変更すると，この確率比の意味が失われる．

mask 付き entropy を
$\mathsf H_n(\theta)=-\sum_{z\in\mathcal Z_n}\pi_\theta(z\mid x_n^{\rm pol},m_n)
\log\pi_\theta(z\mid x_n^{\rm pol},m_n)$ とする．
係数 $c_H,c_V\ge0$ を固定し，最小化する総損失を
\begin{equation}\label{eq:loss}
 \mathcal L(\theta,\nu)=
 -\mathcal L_{\rm clip}(\theta)
 -\frac{c_H}{|I_{\rm act}|}\sum_{n\in I_{\rm act}}\mathsf H_n(\theta)
 +\frac{c_V}{|I_{\rm val}|}\sum_{n\in I_{\rm val}}
       (V_\nu(x_n^{\rm pol})-\hat G_n)^2
\end{equation}
とする．$\hat A_n,\hat G_n$ は更新中は定数として扱う．各適格集合が空ならその項の更新をスキップし，$0$ 除算しない．
GAE の可変時間版 \eqref{eq:gae} は通常の GAE\cite{gae} の割引を各区間の積に置き換えた本稿の提案である．

学習率を $\alpha_{\rm opt}>0$，微分をパラメータごとの偏微分を並べた勾配 $\nabla$ とすると，
最も単純な更新は $(\theta,\nu)\leftarrow(\theta,\nu)-\alpha_{\rm opt}\nabla\mathcal L$ である．
最適化器，隠れ幅，clipping 幅，更新 epoch 数は設定値として固定・保存する．
学習対象が $6$ 戦術だけであることは式 \eqref{eq:policy}--\eqref{eq:loss} 全体で維持する．

\subsection{自己対戦と評価対象}
凍結相手集合を $\mathcal P^-=\{\pi_1^-,\ldots,\pi_K^-\}$，選択重みを
$w_k^-\ge0$，$\sum_{k=1}^K w_k^-=1$ とする．対局開始時に $k$ を抽選し，その対局中は相手を変えない．
この場合の目的は
\begin{equation}\label{eq:opponent-mixture}
 J_{\rm pool}(\theta)=\sum_{k=1}^K w_k^-\,
      \mathbb E_{\mu_0,\pi_k^-,\pi_\theta}[G_0].
\end{equation}
相手を更新したら核 \eqref{eq:kernel} と評価対象分布も変わるので，checkpoint と重みを記録する．
単一相手に勝つことから全戦術・全相手への強さを結論しない．

比較は search-only，ルール戦術，教師初期化，RL 戦術で探索予算と候補生成版を揃える．
評価用の配牌対を $j=1,\ldots,N_e$，左右交換した $2$ 対局の結果を $z_{j,0},z_{j,1}\in\{-1,0,1\}$ とする．
対ごとの勝敗点 $v_j=\tfrac12\sum_{b=0}^1(z_{j,b}+1)/2$ と定義すると，
比較対象との点差 $d_j=v_j^{\rm RL}-v_j^{\rm rule}$ の標本平均は $N_e^{-1}\sum_jd_j$ となる．
$v_j$ は引分を $0.5$ と数える勝敗点であり，純粋な勝利頻度は勝利指示値から別に計算する．
信頼区間の再標本化単位はこの配牌対であり，decision 行や再現確認の repeat を独立標本として水増ししない．

候補不足（公開参照解があるのに候補化できない），戦術内順位失敗（良い候補があるが代表でない），
戦術選択失敗（良い代表戦術があるが選ばない）を別に評価する．
探索で提示できない配置を，戦術だけを学ぶ actor が新しく作ることはできない．
能力検証・短期学習・資源判断・昇格判定は別段階とし，本稿の式の成立だけを能力達成や長時間学習開始の証拠にしない．

\section{実装対応と未実装部分}\label{sec:mapping}
以下のリンクは全て本文の固定 SHA を指す．本文の数学的集合はデータ構造の説明であり，
実数演算がそのまま浮動小数点演算の全ビット一致を保証するわけではない．
\begin{longtable}{p{.31\linewidth}p{.61\linewidth}}
\toprule 本文の対象 & 実装・設計の対応\\\midrule\endhead
盤面，重力，同時消去 & \src{src/core/field.py}{src/core/field.py}：\code{drop_puyo, get_vanish_groups, remove_puyos}\\
配置と得点，全消し & \src{src/core/game.py}{src/core/game.py}：\code{find_landing_y, resolve_chains_synchronously}，得点計算と保留 bit\\
行動 ID & \src{puyo_env/actions.py}{puyo_env/actions.py}：\code{PLACEMENT_ACTIONS}\\
操作の到達可能性 & \src{puyo_env/action_planner.py}{puyo_env/action_planner.py}：\code{plan_placement_action}\\
実時間 tick，勝敗 & \src{src/core/realtime.py}{src/core/realtime.py}，\src{puyo_env/realtime_versus.py}{puyo_env/realtime_versus.py}\\
公開観測，時刻要約 & \src{puyo_env/nextgen_public_snapshot.py}{nextgen_public_snapshot.py}，\src{agents/nextgen_contracts.py}{nextgen_contracts.py}\\
盤面構造評価 & \src{agents/chain_structure.py}{agents/chain_structure.py} と \src{train/config/nextgen_chain_structure.yaml}{nextgen_chain_structure.yaml}\\
シナリオ，beam，根順位 & \src{agents/long_horizon_search.py}{agents/long_horizon_search.py}，\src{agents/nextgen_profiles.py}{nextgen_profiles.py}\\
共有候補，対戦応答 & \src{agents/nextgen_shared_search.py}{nextgen_shared_search.py}，\src{agents/nextgen_response_search.py}{nextgen_response_search.py}\\
公開情報による生存検証 & \src{agents/nextgen_survival.py}{nextgen_survival.py}：\code{apply_envelope}．候補根拠を補強し，戦術の固定順位へ反映\\
テンプレート選択・phase & \src{agents/template_catalog.py}{template_catalog.py}，\src{agents/template_phase.py}{template_phase.py}，\src{train/config/nextgen_templates.yaml}{nextgen_templates.yaml}\\
MVP 戦術選択 & \src{agents/nextgen_tactic_manager.py}{nextgen_tactic_manager.py}：\code{RuleTacticSelector, NextgenTacticManagerPolicy}\\
特徴・mask・実行 receipt & \src{agents/nextgen_contracts.py}{nextgen_contracts.py}：\code{PolicyFeatures, Selection, ExecutionReceipt}\\
学習用記録 & \src{train/nextgen_trajectory.py}{train/nextgen_trajectory.py}：検証・保存・読出し．専用 PPO trainer ではない\\
将来の戦術 RL & \srcja{docs/development/puyo-234-next-generation-model-design.md}{PUYO-234 設計書} 第 13・15・17 節を本稿第 \ref{sec:learning} 節で具体化\\\bottomrule
\end{longtable}

\begin{center}\begin{tabular}{p{.44\linewidth}p{.47\linewidth}}\toprule 実装済みの境界 & 本稿の将来提案\\\midrule
公開 request v2，candidate batch v3，features v1 & 186 入力・6 出力の戦術 actor の実体\\
ルール戦術と固定候補順位 & 教師初期化，時刻差付き PPO/GAE\\
観測，根拠，実行結果を記録する契約 & 次価値を保存・再構成できる rollout collector\\
テンプレート・共有探索・応答探索の固定予算 & 凍結相手 pool による継続学習と採否評価\\\bottomrule\end{tabular}\end{center}
探索予算は shared $600000$，template $128$，response $256$ を別々に数える．
候補探索を戦術数の $6$ 倍だけ繰り返す構成ではない．
定義したモデルの実験には設定，source/native build，feature registry，戦術順，template catalog，
相手 checkpoint，乱数 stream，時刻モードを固定する必要がある．
仕様として定義することと，受け入れ検証に合格することは別である．

\section{読み違いを防ぐ検算例}\label{sec:examples}
\begin{example}[同時相殺]
換算後の生成量が $(g_0,g_1)=(10,8)$，受信残量が $(N_0,N_1)=(4,1)$ なら，
自分の受信との相殺後は $(r_0,r_1)=(6,7)$，同時相殺量は $a_*=6$，
送信量は $(o_0,o_1)=(0,1)$ である．到着前の受信 packet も最初の相殺対象に含める．
\end{example}
\begin{example}[未観測と空]
$\bar B(0,12)=\unknown$ のとき，式 \eqref{eq:completion} はそのセルが空の盤面も赤の盤面も許す．
一方 $\bar B(0,12)=\emptycell$ なら赤の盤面を許さない．
従って未観測セルを $0$ 埋めした単一盤面の生存予測は，公開情報に整合する全盤面での生存証明にはならない．
\end{example}
\begin{example}[mask と log-probability]
6 logit が全て $0$，有効戦術が build\_main と cancel だけなら，式 \eqref{eq:policy} はそれぞれ $1/2$，
ほかを $0$ とする．cancel の保存 log-probability は $-\log2$ である．
更新時に全 $6$ 戦術を有効にし直して $-\log6$ と比較するのは，保存された選択問題と異なる．
\end{example}
\begin{example}[可変時間割引]
区間長が $\Delta_0=0.5$ 秒，$\Delta_1=2$ 秒なら，$\Gamma_0\Gamma_1=\gamma^{2.5}$ となる．
$t_2$ に初めて勝ちが確定し補助なしなら，$R_0^0=0$，$R_1^0=\gamma^2$，
$G_0=\gamma^{2.5}$ である．1 decision ごとに一定の割引を掛けると $\gamma^2$ になり，実時間と一致しない．
\end{example}

\appendix
\section{主要記号の索引}\label{sec:notation-index}
この索引は導入順の定義を参照するためのもので，本文で先に使った記号を事後的に補うものではない．
局所添字と関数引数はそれぞれの式の直前で定義している．
\begin{longtable}{p{.25\linewidth}p{.43\linewidth}p{.22\linewidth}}
\toprule 記号 & 意味 & 定義箇所\\\midrule\endhead
$\N,\R,\Delta(A),\none$ & 非負整数，実数，分布全体，該当値なし & 定義 \ref{def:notation}\\
$\mathcal C,\mathcal V,D,D_v,\mathcal B$ & 色，セル値，全領域，可視領域，盤面集合 & 定義 \ref{def:board}\\
$G,E,\mathcal R,L$ & 重力，消去，連鎖解決，連鎖数 & 式 \eqref{eq:gravity}，\eqref{eq:erase}\\
$p,\mathcal O,\delta,\mathcal A$ & 組ぷよ，向き，子の偏位，配置行動集合 & 式 \eqref{eq:actions}\\
$\mathcal A_{\rm geo},\mathcal A_{\rm reach}$ & 幾何合法配置と操作で到達可能な配置 & 第 \ref{sec:game} 節\\
$s_j,S,b,q,q_{\rm last}$ & 段得点，連鎖得点，保留 bit，総得点，前回連鎖時総得点 & 式 \eqref{eq:score}--\eqref{eq:attackscore}\\
$t,f,\kappa,c_i,g_i$ & tick，tick rate，換算点，繰越点，生成個数 & 式 \eqref{eq:carry}\\
$Q_i,N_i,a_i,a_*,o_i$ & 受信台帳，残量，受信相殺，同時相殺，送信量 & 式 \eqref{eq:cancel}\\
$\mathcal S,F_\eta,\mathcal S_T,z_i$ & 全状態，1 tick 遷移，終端集合，対戦結果 & 定義 \ref{def:state}\\
$\bar B,o_t^i,h_t^i,\xi$ & 部分盤面，公開観測，公開履歴，探索等の独立乱数 & 第 \ref{sec:versus} 節\\
$T,\omega^m,D_s,w,M_s,Q,q_r$ & 単独遷移，シナリオ，深さ，幅，数，予算，根確保枠 & 第 \ref{sec:search} 節\\
$U,\psi,\mathbf w,d_U$ & 固定評価器，特徴，重み，致命判定 & 第 \ref{sec:evaluator} 節\\
$\mathcal L_d^m,\mathcal E_d^m$ & beam の生存列，展開候補列 & 式 \eqref{eq:beam-expand}\\
$c_a,P,\rho_a,K_a$ & 根の発火分類，優先度，coverage，根順位キー & 第 \ref{sec:aggregate} 節\\
$\Sigma,P_T,\beta,\varphi$ & template 記号，形状，色束縛，phase & 第 \ref{sec:tactics} 節\\
$\mathcal Z,\mathcal C_n(z),m_n,c_n^*$ & 戦術集合，候補列，mask，固定代表候補 & 式 \eqref{eq:tactic-mask}\\
$x_n^{\rm pol},y_n,b_n^{\rm miss}$ & actor 入力，raw 要約，欠損 bit 列 & 式 \eqref{eq:actor-input}\\
$\theta,f_\theta,\pi_\theta,\Pi_\theta$ & actor パラメータ，logit，戦術方策，要求配置分布 & 式 \eqref{eq:policy}，\eqref{eq:composed-policy}\\
$\tilde S_n,H_n,b_n$ & 拡張全状態，decision 履歴，信念分布 & 第 \ref{sec:semimarkov} 節\\
$\Delta_n,\Gamma_n,\Lambda_n$ & 区間秒数，割引，trace 減衰 & 式 \eqref{eq:td}--\eqref{eq:gae}\\
$R_n^0,\Phi,\alpha,R_n,G_n,J$ & 基礎報酬，potential，補助係数，報酬，収益，目的 & 式 \eqref{eq:base-reward}--\eqref{eq:return}\\
$V^\pi,Q^\pi,A^\pi,V_\nu$ & 履歴価値，行動価値，advantage，価値近似 & 式 \eqref{eq:value}，\eqref{eq:td}\\
$I_{\rm act},I_{\rm val},\hat A,\hat G$ & actor/value 適格行，GAE，価値教師 & 式 \eqref{eq:gae}--\eqref{eq:loss}\\\bottomrule
\end{longtable}

\section{再現とレビューの手順}\label{sec:reproduce}
文書ソースは \code{docs/mathematical-model/} にある．
同じディレクトリの \code{README.md} に，日本語フォントを含む Tectonic の導入方法とビルドコマンドを記載する．
数式番号は section ごとの LaTeX counter，参照は \code{label/eqref/ref}，出典は \code{cite} で管理する．
番号を本文へ手入力しない．

\code{build.py} はコンパイルと参照検査を行い，エラー・未解決参照・未収録文字・はみ出しを検出したら失敗する．
\code{source-map.json} は対象 SHA と参照ファイルの SHA-256 を保持し，
\code{verify.py} は Git object との一致と検算例を確認する．
自動検査は文章中の数学的正しさや記号の定義順を完全には証明しない．
そのためレビューでは次の順で確認する．
\begin{enumerate}
\item 定義 \ref{def:board} から式 \eqref{eq:placement} までの盤面・配置・消去が同じ座標と状態型を使っているか．
\item 得点・おじゃま・探索評価・学習報酬が区別され，終了と打切りを混同していないか．
\item 式 \eqref{eq:privacy} の境界を候補・特徴・学習サンプルが守っているか．
\item 式 \eqref{eq:composed-policy} の合成から，学習対象が戦術だけであると読めるか．
\item 時刻差と終端 bit が式 \eqref{eq:return}--\eqref{eq:loss} で一貫しているか．
\end{enumerate}

\begin{thebibliography}{9}
\bibitem{repo}
shhchan，puyo-ai-dev-platform，
\href{https://github.com/shhchan/puyo-ai-dev-platform/tree/f8394d77578ecb66e15802fa502d8e1668223006}{固定 commit f8394d77578e}，2026．
\bibitem{design}
PUYO-234，\srcja{docs/development/puyo-234-next-generation-model-design.md}{次世代モデル設計メモ}，同固定 commit 所収．
\bibitem{ppo}
J. Schulman, F. Wolski, P. Dhariwal, A. Radford, O. Klimov,
\emph{Proximal Policy Optimization Algorithms}, 2017.
\url{https://arxiv.org/abs/1707.06347}．
\bibitem{gae}
J. Schulman, P. Moritz, S. Levine, M. Jordan, P. Abbeel,
\emph{High-Dimensional Continuous Control Using Generalized Advantage Estimation}, 2015.
\url{https://arxiv.org/abs/1506.02438}．
\bibitem{tectonic}
The Tectonic Project，\emph{Installing Tectonic}．
\url{https://tectonic-typesetting.github.io/en-US/install.html}．
\end{thebibliography}

\end{document}
